\documentclass{article}
\usepackage[T1]{fontenc}
\usepackage{lmodern}
\usepackage{graphicx}
\usepackage{amsmath,amssymb}
\usepackage[a4paper,margin=2cm]{geometry}
\usepackage[absolute,overlay]{textpos}
\setlength{\TPHorizModule}{1mm}
\setlength{\TPVertModule}{1mm}
\usepackage[indonesian]{babel}
\usepackage{hyperref}
\setlength{\emergencystretch}{2em}
% unit: Halaman 70

\begin{document}

% === Halaman 70 (llm) ===

\begin{itemize}
    \item Algoritma A5/1 telah digunakan untuk mengenkripsi semua percakapan dan komunikasi data melalui telepon seluler GSM di Eropa.
    \item Program A5/1 ditanam di dalam \textit{chip} pada kartu \textit{Simcard}.
    \item Di negara-negara di mana operator telepon seluler dilarang melakukan enkripsi percakapan, seperti di Indonesia, maka program algoritma A5 tidak diaktifkan (\textit{disabled}), sehingga semua sinyal terkirim dalam bentuk plainteks.
\end{itemize}

\end{document}
